\documentclass[12pt, a4paper]{article}
\usepackage{graphicx} % Required for inserting images
\graphicspath{{images/}{images-desafios/}}
\usepackage[a4paper, total={7.5in, 10in}]{geometry}
\usepackage{float} % pra conseguir usar o \begin{figure}[H]
\usepackage{indentfirst} % pra fazer tab n 1 paragrafos dos capitulos
% package pra tabela
\usepackage{multirow}
\usepackage{array}

%pacote pra codigo
\usepackage{listings} 
\usepackage{xcolor}
\definecolor{codegreen}{rgb}{0,0.6,0}
\definecolor{codegray}{rgb}{0.5,0.5,0.5}
\definecolor{codepurple}{rgb}{0.58,0,0.82}
\definecolor{backcolour}{rgb}{0.95,0.95,0.92}

\lstdefinestyle{myCodeStyle}{
	backgroundcolor=\color{backcolour},   
	commentstyle=\color{codegreen},
	keywordstyle=\color{cyan},
	numberstyle=\tiny\color{codegray},
	stringstyle=\color{codepurple},
	basicstyle=\ttfamily\footnotesize,
	breakatwhitespace=false,         
	breaklines=true,                 
	captionpos=b,                    
	keepspaces=true,                 
	%numbers=left,                    
	numbersep=5pt,                  
	showspaces=false,                
	showstringspaces=false,
	showtabs=false,                  
	tabsize=2
}
\lstset{style=myCodeStyle}
\renewcommand{\lstlistingname}{Código}%Muda o termo "Listing":  Listing -> Código
\renewcommand{\lstlistlistingname}{Lista de \lstlistingname s}% List of Listings -> Lista de Códigos


\newcommand{\unit}[1]{\ensuremath{\, \mathrm{#1}}} % comando pra tirar do math mode a unidade de medida

\def\tamanhoGrafico{0.8}

\newcommand{\ohm}{\unit{\Omega} }
\newcommand{\kohm}{\unit{k\Omega}}

\title{
	Laboratório Digital A - PCS3335
	\\ Relatório da Experiência 6
}

\author{
	Alunos:
	\\ Rafael Hipólito Rodrigues - N.USP: 14604308
	\\ {{PERSON-2}} - N.USP: {{PERSON-ID-2}}
	\\ \\
	Turma 3 - Bancada A6
	\\ Professor: Bruno de Carvalho Albertini
}

%\date{}

\begin{document}
	
	\maketitle
	
	\section*{Planejamento} 
	Elaboração do circuito de recebimento de dados(bits) numa UART simplificada.
	O circuito responsável pelo recebimento é a entidade \texttt{Receiver}.
	Foi mantido as entidades de transmissão da experiência 5.
	
	\section*{Diagramas}
	
	\subsection*{FSM}
	
	Diagrama da Maquina de Estados Finitos do Receiver gerado no Quartus : 
	
	\begin{figure}[H]
		\centering
		\includegraphics[width=0.9\linewidth]{images/FSM.png}
		\caption{FSM do Receiver.}
		\label{fig:FSM_exp4}
	\end{figure}
	
	
	
	\subsection*{RTL}
	
	Diagramas RTL da entidades gerado no Quartus:
	
	\begin{figure}[H]
		\centering
		\includegraphics[width=0.9\linewidth]{images/RTL_uart.png}
		\caption{Diagrama RTL da UART.}
		\label{fig:RTL1}
	\end{figure}
	
	\begin{figure}[H]
		\centering
		\includegraphics[width=0.9\linewidth]{images/RTL_receiver.png}
		\caption{Diagrama RTL do \texttt{Receiver}.}
		\label{fig:RTL_TCC}
	\end{figure}

	
	Os contadores e registradores utilizados foram os mesmos da experiência 2.
	
	\section*{Pinagem}

	
	\begin{table}[H]
		\begin{center}
			\begin{tabular}{|l|l|}
				\hline
				busLine{[}7{]}     & PIN\_AA13 \\ \hline
				busLine{[}6{]}     & PIN\_AA14 \\ \hline
				busLine{[}5{]}     & PIN\_AB15 \\ \hline
				busLine{[}4{]}     & PIN\_AA15 \\ \hline
				busLine{[}3{]}     & PIN\_T12  \\ \hline
				busLine{[}2{]}     & PIN\_T13  \\ \hline
				busLine{[}1{]}     & PIN\_V13  \\ \hline
				busLine{[}0{]}     & PIN\_U13  \\ \hline
				busSelect{[}0{]}   & PIN\_AB13 \\ \hline
				clock              & PIN\_M9   \\ \hline
				lsrDebugger{[}6{]} & PIN\_L1   \\ \hline
				lsrDebugger{[}5{]} & PIN\_L2   \\ \hline
				reset              & PIN\_B16  \\ \hline
				serial\_out        & PIN\_n16  \\ \hline
			\end{tabular}
		\end{center}
	\end{table}
	
	
	
	\section*{Testes}
	
	\subsection*{Resultados dos testes}

	\begin{itemize}
		\item bits adequadamente envidados em 9600bps.
		\par Resultado: SUCESSO
		
		\item Character Length
		\begin{table}[H]
			\begin{center}
				\begin{tabular}{|l|l|l|l|}
					\hline
					LCR\_bit 0 & LCR\_bit 1 & Character Length & Resultado \\ \hline
					0     & 0     & 5 bits           &SUCESSO           \\ \hline
					0     & 1     & 6 bits           &SUCESSO           \\ \hline
					1     & 0     & 7 bits           &SUCESSO           \\ \hline
					1     & 1     & 8 bits           &SUCESSO           \\ \hline
				\end{tabular}
			\end{center}
		\end{table}
		
		\item Bit Stop
		\begin{table}[H]
			\begin{center}
				\begin{tabular}{|l|l|l|l|l|}
					\hline
					LCR\_bit 2 & LCR\_bit 0 & LCR\_bit 1 & Quant. bitStop & Resultado\\ \hline
					0          & \#         & \#         & 1              & SUCESSO  \\ \hline
					1          & 0          & 0          & 1,5            & SUCESSO  \\ \hline
					1          & 0          & 1          & 2              & SUCESSO  \\ \hline
					1          & 1          & 0          & 2              & SUCESSO  \\ \hline
					1          & 1          & 1          & 2              &  SUCESSO \\ \hline
				\end{tabular}
			\end{center}
		\end{table}
		
		\item Bit Paridade
		\begin{table}[H]
			\begin{center}
				\begin{tabular}{|l|l|l|l|l|}
					\hline
					LCR\_bit 3 & LCR\_bit 5 & LCR\_bit 4 & bit Paridade     & Resultado \\ \hline
					0          & \#         & \#         & SEM bit paridade & SUCESSO  \\ \hline
					1          & 0          & 0          & Paridade IMPAR   & SUCESSO   \\ \hline
					1          & 0          & 1          & Paridade PAR     & SUCESSO   \\ \hline
					1          & 1          & 0          & 1                & SUCESSO   \\ \hline
					1          & 1          & 1          & 0                & SUCESSO   \\ \hline
				\end{tabular}
			\end{center}
		\end{table}
		
		
		\item Ao ativar e desativa o reset, volta ao estado inicial. 
		\par Resultado: SUCESSO
		
	\end{itemize}
	
	
\end{document}
